\chapter{Le bureau : Windows, Linux, macOS}

Trois systèmes, trois façons de compiler, et un descripteur unique. Ce chapitre
donne ce que chacun exige --- et ce qu'il exige \emph{de la machine}, qui n'est
pas la même chose.

\section{Les chaînes d'outils disponibles}

\begin{nkcode}{jenga info -- capture reelle, extrait}
Name                 Family        Target OS   Arch     Env
===============================================================
host-clang           clang         Windows     x86_64   mingw
host-gcc             gcc           Windows     x86_64   mingw
msvc                 msvc          Windows     x86_64   msvc
clang-mingw          clang         Windows     x86_64   mingw
mingw                gcc           Windows     x86_64   mingw
clang-cross-linux    clang         Linux       x86_64   gnu
android-ndk          android-ndk   Android     arm64    android
emscripten           emscripten    Web         wasm32
\end{nkcode}

\begin{nkretenir}
\textbf{Cette table est celle de LA MACHINE, pas celle de Jenga.} Android et
Emscripten y figurent parce que leurs SDK sont installés ici. Sur une machine
neuve, la table est plus courte --- et un descripteur qui exige une chaîne
absente échoue au chargement.

Jenga en déclare onze au total ; la table montre celles qu'il a \emph{trouvées}.
\end{nkretenir}

\section{Windows}

\begin{nkcode}{Applications/NkDames/NkDames.jenga:113-129 -- cite tel quel}
with filter("system:Windows && !options:windows-runtime=uwp && !system:XboxSeries && !system:XboxOne"):
    windowedapp()
    usetoolchain(TC_WINDOWS)
    defines(["WIN32_LEAN_AND_MEAN", "_UNICODE", "UNICODE", _VK_ON])
    if _VULKAN_LIB:
        libdirs([_VULKAN_LIB])
    _WIN_LINKS = [
        "user32", "gdi32", "opengl32", "dwmapi", "shell32",
        "mf", "mfplat", "mfreadwrite", "mfuuid",
        "uuid", "ole32", "dinput8", "dxguid",
        "d3d11", "d3d12", "dxgi", "d3dcompiler",
        "avrt", "winmm",   # NKAudio (WASAPI / MMCSS)
    ]
    if _WANT_VULKAN:
        _WIN_LINKS.append("vulkan-1")
    links(_WIN_LINKS)
\end{nkcode}

\begin{nkretenir}[Ce bloc est la thèse du livre en huit lignes]
\textbf{Rien ici n'est une « fonctionnalité de Jenga ».} La liste est construite
par du Python ordinaire : on l'écrit, on lui ajoute conditionnellement une
entrée, puis on la \emph{passe} à \texttt{links()}. C'est la conséquence directe
du chapitre~\ref{ch:python} --- un \texttt{.jenga} est exécuté.

Un système de construction à syntaxe déclarative demanderait ici une notion
spéciale de « liste conditionnelle ». Il n'y en a pas besoin : \texttt{if} et
\texttt{append} existaient déjà.
\end{nkretenir}

\begin{nkpiege}[\texttt{TC\_WINDOWS} n'est pas de Jenga]
\alerte{}\textbf{Ne le tapez pas tel quel : vous obtiendrez un
\texttt{NameError}.} \texttt{usetoolchain} prend une \emph{chaîne} ---
\texttt{def usetoolchain(name: str)}, \nkfile{Jenga/Core/Api.py:1836} --- et
les exemples livrés avec Jenga écrivent \texttt{usetoolchain("clang-mingw")}.

\texttt{TC\_WINDOWS} est une variable de \textbf{notre} dépôt :

\begin{nkcode}{config/toolchain.jenga:14}
TC_WINDOWS = "nk-windows-clang-mingw"
\end{nkcode}

Un \texttt{.jenga} étant du Python, une variable définie dans un fichier inclus
est visible partout ensuite --- c'est ce qui permet de nommer une fois la chaîne
d'outils et de la changer pour \textbf{189 descripteurs} d'un coup. Le prix est
qu'à la lecture, rien ne distingue un symbole de Jenga d'un symbole maison.

\textbf{C'est le deuxième de la série}, après \texttt{nkentseudependson}
(chapitre~\ref{ch:dependances}). Devant un nom que la documentation de Jenga
ignore, la question est toujours la même : \emph{où est-il défini ?} Un
\texttt{grep} sur \texttt{<nom> =} dans vos propres \texttt{.jenga} répond en
une seconde.
\end{nkpiege}

\begin{nkretenir}
\textbf{\texttt{WIN32\_LEAN\_AND\_MEAN} n'est pas de la superstition.} Sans lui,
\texttt{windows.h} tire des dizaines de milliers de lignes dont vous n'avez pas
besoin --- et surtout des \textbf{macros} : \texttt{min}, \texttt{max},
\texttt{near}, \texttt{far}, \texttt{small}. Chacune casse silencieusement un
code qui emploie le même nom.

\textbf{\texttt{\_UNICODE} et \texttt{UNICODE}} choisissent les variantes larges
des API. Les deux, parce que l'un est lu par la bibliothèque C et l'autre par
l'API Windows --- en déclarer un seul donne un mélange qui compile et se comporte
mal sur les chemins accentués.
\end{nkretenir}

\begin{nkretenir}[Les bibliothèques se déclarent par leur nom nu]
\texttt{"user32"} et non \texttt{"user32.lib"}. Jenga ajoute ce qu'il faut selon
la chaîne : \texttt{-luser32} pour MinGW, \texttt{user32.lib} pour MSVC. C'est ce
qui rend la ligne identique pour les deux.
\end{nkretenir}

\section{Linux, et ses quatre backends}

C'est le cas le plus instructif du livre, parce qu'un même système y demande
quatre descriptions.

\begin{nkcode}{Applications/NkDames/NkDames.jenga -- reel, abrege}
with filter("system:Linux && options:linux-backend=xlib || "
            "system:Linux && !options:linux-backend && !options:headless"):
    usetoolchain("clang-native")
    objdir(".../Build/Obj/%{cfg.buildcfg}-%{cfg.system}-xlib/%{prj.name}")
    targetdir(".../Build/Bin/%{cfg.buildcfg}-%{cfg.system}-xlib/%{prj.name}")
    defines(["NKENTSEU_FORCE_WINDOWING_XLIB_ONLY"])
    links(["pthread", "X11", "Xext", "Xrandr", "GL", "asound"])
\end{nkcode}

\begin{center}
\begin{tabular}{@{}lll@{}}
\toprule
\textbf{Backend} & \textbf{Ce qu'il lie} & \textbf{Sortie} \\
\midrule
\texttt{xlib} & \texttt{X11 Xext Xrandr GL asound} & \texttt{Release-Linux-xlib} \\
\texttt{xcb} & \texttt{xcb xcb-image xkbcommon\dots} & \texttt{Release-Linux-xcb} \\
\texttt{wayland} & \texttt{wayland-client wayland-egl\dots} & \texttt{Release-Linux-wayland} \\
\texttt{headless} & \texttt{pthread} seulement & \texttt{Release-Linux-headless} \\
\bottomrule
\end{tabular}
\end{center}

\begin{nkretenir}
\textbf{Chaque backend a son propre dossier de sortie}, et c'est indispensable :
sans le suffixe, construire pour Wayland écraserait la construction XLib, et
l'on lancerait un binaire lié aux mauvaises bibliothèques.

\alerte{} \textbf{Le corollaire coûte cher aux scripts} : le dossier n'est pas
\texttt{Release-Linux} mais \texttt{Release-Linux-xlib}. Un empaqueteur qui code
le premier en dur échoue sur une construction \textbf{réussie} --- mesuré sur ce
dépôt.
\end{nkretenir}

\begin{nkretenir}[Le second membre de l'expression est le défaut]
Relisez la condition : « backend xlib demandé, \emph{ou} aucun backend demandé et
pas de mode sans écran ». Le second membre attrape l'utilisateur qui n'a rien
choisi.

C'est la façon d'écrire un défaut avec des filtres --- il n'y a pas de mot-clé
« sinon ».
\end{nkretenir}

\begin{nkpiege}
\textbf{Construire pour Linux depuis Windows ne marche pas} avec cette
description : \texttt{clang-native} y cherche \texttt{X11/X.h}, qui n'existe pas.

Deux issues, et le dépôt emploie les deux : \textbf{WSL2}, où les en-têtes
existent, pour les constructions locales ; une \textbf{intégration continue}
Linux pour le reste. Aucune n'est un contournement : ce sont des machines Linux.
\end{nkpiege}

\section{macOS}

\begin{nkcode}{Applications/NkDames/NkDames.jenga -- reel}
with filter("system:macOS"):
    windowedapp()
    usetoolchain("clang-native")
    frameworks(["Cocoa", "QuartzCore", "OpenGL", "Metal",
                "AudioToolbox", "CoreAudio", "AudioUnit"])
\end{nkcode}

\begin{nkretenir}
\textbf{Pas de \texttt{links}, des \texttt{frameworks}.} Un framework Apple porte
à la fois les en-têtes et le binaire.

Le commentaire du descripteur note que cette liste a été établie \emph{sur
l'édition de liens de l'intégration continue}, pas devinée --- et il précise le
point qui surprend : « les 23 projets COMPILENT sans ces frameworks ; seul le
lien final tombe ».

\emph{Une dépendance de framework ne se voit pas à la compilation.} C'est
pourquoi elle s'oublie, et pourquoi elle ne se découvre qu'à la toute fin.
\end{nkretenir}

\begin{nkpiege}
\textbf{Construire pour macOS exige macOS} --- et Xcode. Jenga déclare des
chaînes de compilation croisée, mais le SDK Apple ne se redistribue pas.

Le dépôt passe donc par une intégration continue macOS, et son script
d'empaquetage le dit en toutes lettres plutôt que de laisser croire le
contraire.
\end{nkpiege}

\section{Debug et Release}

\begin{nkcode}{La forme, identique dans tous les descripteurs}
with filter("config:Debug"):
    defines(["_DEBUG", "DEBUG"]); optimize("Off");   symbols(True)
with filter("config:Release"):
    defines(["NDEBUG"]);          optimize("Speed"); symbols(False)
\end{nkcode}

\begin{nkretenir}
\textbf{Debug et Release ne sont pas le même programme}, et il faut le prendre au
sérieux : les \texttt{assert} disparaissent, la mémoire non initialisée change de
contenu, et l'optimiseur expose des comportements indéfinis que le Debug
masquait.

\textbf{Un défaut qui n'existe que dans l'une des deux configurations est
fréquent.} Toute mesure, tout chiffre, toute campagne doit donc porter sa
configuration --- sans quoi deux mesures justes se contredisent sans que
personne ait tort.
\end{nkretenir}

\begin{nkbilan}
Windows demande \texttt{WIN32\_LEAN\_AND\_MEAN} et les deux macros Unicode, et
ses bibliothèques se nomment nues. Linux n'est pas une cible mais quatre, chacune
avec son dossier de sortie suffixé --- ce que les scripts doivent lire et non
coder en dur --- et le second membre du filtre y sert de valeur par défaut. macOS
emploie des frameworks, dont la liste ne se découvre qu'à l'édition de liens.
Enfin Debug et Release produisent deux programmes différents : une mesure sans sa
configuration n'est pas une mesure.
\end{nkbilan}

\begin{nkquestions}
\begin{enumerate}
\item La table des chaînes d'outils décrit-elle Jenga ou la machine ?
\item Que se passe-t-il sans \texttt{WIN32\_LEAN\_AND\_MEAN} ?
\item Pourquoi chaque backend Linux a-t-il son propre dossier de sortie ?
\item Pourquoi un framework macOS manquant ne se voit-il qu'à la fin ?
\item Citez trois différences réelles entre Debug et Release.
\end{enumerate}
\end{nkquestions}

\begin{nkpratique}
Faites construire votre projet du chapitre 3 sur deux configurations et vérifiez,
\emph{sur le disque}, que les deux artefacts coexistent.

Puis ajoutez un filtre Linux avec une option \texttt{-{}-mon-backend} à trois
valeurs, chacune posant une définition différente et un dossier de sortie
différent.

Vérifiez par \texttt{jenga compile-flags} --- pas par la construction, que vous
ne pouvez peut-être pas faire ici. \textbf{C'est justement l'intérêt} : on peut
vérifier une description pour une plateforme qu'on n'a pas.
\end{nkpratique}

\begin{nkdemo}
Prouvez que Debug et Release ne produisent pas le même programme.

Écrivez un programme dont le comportement dépend de \texttt{NDEBUG} --- un
\texttt{assert} qui échoue, par exemple, ou un \texttt{\#ifdef} qui change un
message. Construisez et lancez dans les deux configurations.

\textbf{Ce que la démonstration doit établir} : ce ne sont pas deux vitesses du
même programme, ce sont deux programmes.

Puis, plus utile encore : \textbf{mesurez la taille des deux exécutables}. L'écart
vous donnera une idée de ce que les symboles de débogage coûtent --- et vous
saurez, la prochaine fois qu'un binaire vous paraît trop gros, quelle question
poser en premier.
\end{nkdemo}

\begin{nklogique}
Un collègue rapporte : « le programme plante en Release, mais pas en Debug ».

\begin{enumerate}
\item Énumérez trois familles de causes possibles, propres à cette asymétrie.
\item Pourquoi le réflexe « je vais déboguer ça en Debug » est-il précisément
      celui qui ne trouvera rien ?
\item Proposez deux approches qui gardent le Release --- et dites ce que chacune
      coûte.
\end{enumerate}
\end{nklogique}
